\section{课程框架与技术背景}

\begin{overviewbox}
课程运行框架包括课程目标、资源平台、作业流程、评分方式、计算环境以及学术诚信要求。单核频率不再持续快速增长之后，现代计算系统转向\term{并行性}{parallelism}、\term{异构性}{heterogeneity}与\term{处理器专用化}{processor specialization}，图形处理器也成为这一转向中的核心加速器。
\end{overviewbox}

相关硬件趋势与软件责任构成以下因果链：
\begin{equation*}
\begin{aligned}
&\text{晶体管密度持续增加}
\quad + \quad
\text{频率缩放趋于停滞}
\\
&\Longrightarrow
\text{更多晶体管无法自动转化为更快的单线程执行}
\\
&\Longrightarrow
\text{多核、向量化与专用加速器成为主要性能路径}
\\
&\Longrightarrow
\text{软件必须显式暴露并行性，并控制负载、访存与串行瓶颈}.
\end{aligned}
\end{equation*}

\subsection{课程的技术定位}

课程名称为应用并行编程，核心对象是\term{大规模并行处理器}{massively parallel processor}，主要以支持统一计算设备架构的英伟达图形处理器为实验平台。课程并不把处理器结构本身作为最终目标；架构知识只学习到足以解释程序性能的深度。换言之，处理器结构、编程接口、并行算法模式和性能测量在本课程中始终服务于同一个目标：把一个具有足够并行性的计算映射到实际硬件上，并使映射后的程序获得可解释、可复现且可扩展的性能。

课程目标包含以下五类技术内容：
\begin{itemize}
  \item 并行编程的基本模型与术语；
  \item 常见并行算法的原则与模式；
  \item 编程接口、工具链与实现技巧；
  \item 影响性能的处理器结构特征与约束；
  \item 以人工智能负载为代表的高价值应用。
\end{itemize}

\begin{keybox}
并行程序的质量不能只用“能否运行”衡量。一个合格的并行实现至少同时满足四项要求：结果正确、性能足够高、代码可维护、算法与实现能适应后续硬件代际。仅通过增加线程数得到一次性的速度提升，并不等于获得了可扩展的并行程序。
\end{keybox}

\subsection{课程形成背景与计算平台脉络}

美国国家超级计算应用中心的图形处理器系统建设构成了本课程的计算平台背景。课程教师长期从事高性能计算、计算加速器以及机器学习系统研究。三代具有代表性的系统分别是：2008 年建设的加速集群使用较早期的图形处理器；2018 年建设的面向人工智能的集群使用 V100 图形处理器；2024 年建设的 DeltaAI 系统采用基于 GH200 超级芯片的 Grace--Hopper 节点。课程实验实际使用较早的 Delta 超级计算机，而不是 DeltaAI。

图形处理器属于更大的系统演化趋势：高性能计算系统逐步从以通用中央处理器为绝对中心的结构，转向由中央处理器负责控制与串行工作、由一个或多个加速器承担大规模并行计算的异构结构。

\subsection{课程规则与技术导论的逻辑关系}

课程运行规则与技术导论看似是两组独立内容，实际存在直接联系：
\begin{itemize}
  \item 课程要求在 Delta 上编译并通过作业调度系统运行，是因为真实超级计算机由大量用户共享，计算资源必须被隔离和调度。
  \item 实验强调自己编写代码，是因为期中考试要求在纸面上完成与实验相近的程序片段；依赖自动代码生成无法形成独立实现能力。
  \item 项目同时评价功能与性能，是因为并行程序的价值来自正确性与性能的共同满足。
  \item 硬件代际快速变化，因此并行程序需要具备可移植性、可扩展性和可维护性，不能只适用于某一张图形处理器。
\end{itemize}

评估并行化方法时需要同时回答两个问题：第一，该方法为什么能提高吞吐量；第二，该方法在真实共享系统、未来硬件和完整应用中是否仍然成立。
